\chapter{2001年豫粤卷}
\section{单选题}
\begin{problem}
不等式 \(\frac{x - 1}{x - 3} > 0\) 的解集为（\quad）

\choices
    {\(\{x \mid x < 1\}\)}
    {\(\{x \mid x > 3\}\)}
    {\(\{x \mid x < 1 \text{或} x > 3\}\)}
    {\(\{x \mid 1 < x < 3\}\)}
\end{problem}

\begin{answer}
C
\end{answer}

\begin{solution}
不等式的定义域为 \(x\ne3\). 分子和分母的零点分别为 \(x=1\) 和 \(x=3\)，将数轴分为

\((-\infty,1)\)，\((1,3)\)，\((3,+\infty)\).

当 \(x<1\) 时，\(x-1<0\)，\(x-3<0\)，分式为正；当 \(1<x<3\) 时，分子为正、分母为负，分式为负；当 \(x>3\) 时，分子和分母都为正，分式为正. \(x=1\) 不能使严格不等式成立，\(x=3\) 不在定义域内，所以解集为 \((-\infty,1)\cup(3,+\infty)\)，故选 C.
\end{solution}

\begin{problem}
若一个圆锥的轴截面是等边三角形，其面积为 \(\sqrt{3}\)，则这个圆锥的全面积是（\quad）

\choices
    {\(3\pi\)}
    {\(3\sqrt{3}\pi\)}
    {\(6\pi\)}
    {\(9\pi\)}
\end{problem}

\begin{answer}
A
\end{answer}

\begin{solution}
设圆锥底面半径为 \(r\)，母线长为 \(l\). 圆锥的轴截面以底面直径 \(2r\) 为一边，且轴截面是等边三角形，所以其三边均为 \(2r\)，从而 \(l=2r\). 由轴截面面积为 \(\sqrt{3}\)，得

\[
\frac{\sqrt{3}}{4}(2r)^2=\sqrt{3}r^2=\sqrt{3}
\]

所以 \(r=1\)，\(l=2\). 因此圆锥的全面积为

\[
\pi r^2+\pi rl=\pi\cdot1^2+\pi\cdot1\cdot2=3\pi
\]

故选 A.
\end{solution}

\begin{problem}
极坐标方程 \(\rho^2\cos 2\theta = 1\) 所表示的曲线是（\quad）

\choices
    {两条相交直线}
    {圆}
    {椭圆}
    {双曲线}
\end{problem}

\begin{answer}
D
\end{answer}

\begin{solution}
由极坐标与直角坐标的关系 \(x=\rho\cos\theta\)，\(y=\rho\sin\theta\)，有

\[
\rho^2\cos2\theta=\rho^2\left(\cos^2\theta-\sin^2\theta\right)=x^2-y^2
\]

所以原方程等价于 \(x^2-y^2=1\)，这是以坐标轴为渐近线的一条双曲线，故选 D.
\end{solution}

\begin{problem}
若定义在区间 \((-1,0)\) 内的函数 \(f(x) = \log_{2a}(x + 1)\) 满足 \(f(x) > 0\)，则 \(a\) 的取值范围是（\quad）

\choices
    {\(\left(0, \frac{1}{2}\right)\)}
    {\(\left(0, \frac{1}{2}\right]\)}
    {\(\left(\frac{1}{2}, +\infty\right)\)}
    {\((0, +\infty)\)}
\end{problem}

\begin{answer}
A
\end{answer}

\begin{solution}
由于 \(f(x)=\log_{2a}(x+1)\) 在所给区间有定义，必须有 \(2a>0\) 且 \(2a\ne1\). 对任意 \(x\in(-1,0)\)，都有 \(0<x+1<1\)，所以 \(\ln(x+1)<0\). 当底数 \(0<2a<1\) 时，对数函数在正数范围内为减函数，故

\[
\log_{2a}(x+1)>\log_{2a}1=0
\]

若 \(2a>1\)，则 \(\log_{2a}(x+1)<0\)，不符合题意. 因而 \(0<2a<1\)，即 \(0<a<\frac12\)，故选 A.
\end{solution}

\begin{problem}
已知复数 \(z = \sqrt{2} + \sqrt{6}\i\)，则 \(\arg \frac{1}{z}\) 是（\quad）

\choices
    {\(\frac{\pi}{3}\)}
    {\(\frac{5\pi}{3}\)}
    {\(\frac{\pi}{6}\)}
    {\(\frac{11\pi}{6}\)}
\end{problem}

\begin{answer}
B
\end{answer}

\begin{solution}
先求倒数：

\[
\frac{1}{z}=\frac{\sqrt{2}-\sqrt{6}\i}{(\sqrt{2})^2+(\sqrt{6})^2}=\frac{\sqrt{2}-\sqrt{6}\i}{8}
\]

该复数的实部为正、虚部为负，位于第四象限，且其辐角的锐角满足

\(\tan\theta=\frac{\sqrt{6}}{\sqrt{2}}=\sqrt{3}\)，\(\theta=\frac{\pi}{3}\).

因此在 \(0\leqslant\arg z<2\pi\) 的约定下，

\[
\arg\frac{1}{z}=2\pi-\frac{\pi}{3}=\frac{5\pi}{3}
\]

故选 B.
\end{solution}

\begin{problem}
函数 \(y = 2^{-x} + 1\)（\(x > 0\)）的反函数是（\quad）

\choices
    {\(y = \log_2 \frac{1}{x - 1}\)（\(x \in (1, 2)\)）}
    {\(y = -\log_2 \frac{1}{x - 1}\)（\(x \in (1, 2)\)）}
    {\(y = \log_2 \frac{1}{x - 1}\)（\(x \in (1, 2]\)）}
    {\(y = -\log_2 \frac{1}{x - 1}\)（\(x \in (1, 2]\)）}
\end{problem}

\begin{answer}
A
\end{answer}

\begin{solution}
由 \(y=2^{-x}+1\) 得 \(y-1=2^{-x}\)，所以

\(-x=\log_2(y-1)\)，\(y=-\log_2(y-1)=\log_2\frac{1}{y-1}\).

原函数的定义域为 \(x>0\)，且 \(0<2^{-x}<1\)，所以其值域为 \(1<y<2\). 交换自变量与因变量的记号后，反函数为

\(y=\log_2\frac{1}{x-1}\)，\(x\in(1,2)\).

故选 A.
\end{solution}

\begin{problem}
若 \(0 < \alpha < \beta < \frac{\pi}{4}\)，\(\sin \alpha + \cos \alpha = a\)，\(\sin \beta + \cos \beta = b\)，则（\quad）

\choices
    {\(a > b\)}
    {\(a < b\)}
    {\(ab < 1\)}
    {\(ab > 2\)}
\end{problem}

\begin{answer}
B
\end{answer}

\begin{solution}
设 \(h(x)=\sin x+\cos x\). 在 \(0<x<\frac{\pi}{4}\) 时，

\[
h'(x)=\cos x-\sin x>0
\]

所以 \(h(x)\) 在该区间内单调递增. 由 \(0<\alpha<\beta<\frac{\pi}{4}\)，得

\[
a=h(\alpha)<h(\beta)=b
\]

故选 B.
\end{solution}

\begin{problem}
在正三棱柱 \(ABC - A_{1}B_{1}C_{1}\) 中，若 \(AB = \sqrt{2}BB_{1}\)，则 \(AB_{1}\) 与 \(C_{1}B\) 所成的角的大小为（\quad）

\choices
    {\(60^{\circ}\)}
    {\(90^{\circ}\)}
    {\(45^{\circ}\)}
    {\(120^{\circ}\)}
\end{problem}

\begin{answer}
B
\end{answer}

\begin{solution}
设正三棱柱的底面边长为 \(s\)，侧棱长为 \(h\)，则题设给出 \(s=\sqrt{2}h\). 建立直角坐标系，取

\(A(0,0,0)\)，\(B(s,0,0)\)，\(C\left(\frac{s}{2},\frac{\sqrt{3}s}{2},0\right)\)

并令 \(A_1\)，\(B_1\)，\(C_1\) 的竖坐标均为 \(h\). 则

\(\overrightarrow{AB_1}=(s,0,h)\)，\(\overrightarrow{C_1B}=\left(\frac{s}{2},-\frac{\sqrt{3}s}{2},-h\right)\).

两向量的数量积为

\[
\overrightarrow{AB_1}\cdot\overrightarrow{C_1B}
=\frac{s^2}{2}-h^2=\frac{2h^2}{2}-h^2=0
\]

所以直线 \(AB_1\) 与 \(C_1B\) 垂直，所成的角为 \(90^{\circ}\)，故选 B.
\end{solution}

\begin{problem}
设 \(f(x)\)，\(g(x)\) 都是单调函数，有如下四个命题：

\begin{circlelist}
\item 若 \(f(x)\) 单调递增，\(g(x)\) 单调递增，则 \(f(x) - g(x)\) 单调递增；
\item 若 \(f(x)\) 单调递增，\(g(x)\) 单调递减，则 \(f(x) - g(x)\) 单调递增；
\item 若 \(f(x)\) 单调递减，\(g(x)\) 单调递增，则 \(f(x) - g(x)\) 单调递减；
\item 若 \(f(x)\) 单调递减，\(g(x)\) 单调递减，则 \(f(x) - g(x)\) 单调递减.
\end{circlelist}

其中正确的命题是（\quad）

\choices
    {\circled{1}\circled{3}}
    {\circled{1}\circled{4}}
    {\circled{2}\circled{3}}
    {\circled{2}\circled{4}}
\end{problem}

\begin{answer}
C
\end{answer}

\begin{solution}
任取 \(x_1<x_2\)，记 \(F(x)=f(x)-g(x)\).

若 \(f\)，\(g\) 都单调递增，取 \(f(x)=x\)，\(g(x)=2x\)，则 \(F(x)=-x\) 单调递减，命题\(\circled{1}\)不成立.

若 \(f\) 单调递增、\(g\) 单调递减，则

\[
F(x_2)-F(x_1)=\bigl(f(x_2)-f(x_1)\bigr)+\bigl(g(x_1)-g(x_2)\bigr)\geqslant0
\]

所以 \(F\) 单调递增，命题\(\circled{2}\)成立.

若 \(f\) 单调递减、\(g\) 单调递增，则

\[
F(x_2)-F(x_1)=\bigl(f(x_2)-f(x_1)\bigr)-\bigl(g(x_2)-g(x_1)\bigr)\leqslant0
\]

所以 \(F\) 单调递减，命题\(\circled{3}\)成立.

若 \(f\)，\(g\) 都单调递减，取 \(f(x)=-x\)，\(g(x)=-2x\)，则 \(F(x)=x\) 单调递增，命题\(\circled{4}\)不成立. 因此正确的命题是\(\circled{2}\)、\(\circled{3}\)，故选 C.
\end{solution}

\begin{problem}
对于抛物线 \(y^{2} = 4x\) 上任意一点 \(Q\)，点 \(P(a, 0)\) 都满足 \(|PQ| \geqslant |a|\)，则 \(a\) 的取值范围是（\quad）

\choices
    {\((-\infty, 0)\)}
    {\((-\infty, 2]\)}
    {\([0, 2]\)}
    {\((0, 2)\)}
\end{problem}

\begin{answer}
B
\end{answer}

\begin{solution}
将抛物线 \(y^2=4x\) 上的点参数化为

\(Q(t^2,2t)\)，\(t\in\R\).

于是

\[
|PQ|^2=(t^2-a)^2+(2t)^2=a^2+t^2\left(t^2+4-2a\right)
\]

题设要求对任意 \(t\in\R\) 都有 \(|PQ|^2\geqslant a^2\). 若 \(a\leqslant2\)，则 \(t^2+4-2a\geqslant0\)，条件成立. 若 \(a>2\)，取满足 \(0<t^2<2a-4\) 的 \(t\)，则 \(t^2+4-2a<0\)，从而 \(|PQ|^2<a^2\)，与题设矛盾. 所以 \(a\leqslant2\)，故选 B.
\end{solution}

\begin{problem}
一间民房的屋顶有如图三种不同的盖法，分别是：

\begin{circlelist}
\item 单向倾斜；
\item 双向倾斜；
\item 四向倾斜.
\end{circlelist}

记三种盖法屋顶面积分别为 \(P_{1}\)，\(P_{2}\)，\(P_{3}\). 若屋顶斜面与水平面所成的角都是 \(\alpha\)，则（\quad）

\problemresetlayout
\begingroup
\expandafter\def\csname bitmaplayout@img/2001/henan_guangdong/q11_fig1.png\endcsname{trim={0bp 0bp 0bp 0bp},clip,width=7.7cm,max width=\linewidth}
\bitmapfigure{img/2001/henan_guangdong/q11_fig1.png}
\endgroup

\choices
    {\(P_{3} > P_{2} > P_{1}\)}
    {\(P_{3} > P_{2} = P_{1}\)}
    {\(P_{3} = P_{2} > P_{1}\)}
    {\(P_{3} = P_{2} = P_{1}\)}
\end{problem}

\begin{answer}
D
\end{answer}

\begin{solution}
设三种盖法的屋顶斜面总面积分别为 \(P_1\)，\(P_2\)，\(P_3\)，房屋水平投影的面积记为 \(S\). 一个斜面与水平面成角 \(\alpha\) 时，该斜面在水平面上的投影面积等于其面积乘以 \(\cos\alpha\). 每一种盖法的各个斜面在水平投影上恰好拼成同一个房屋水平投影，因此都有

\(P_1\cos\alpha=S\)，\(P_2\cos\alpha=S\)，\(P_3\cos\alpha=S\).

由于三种盖法的 \(\alpha\) 相同，得 \(P_1=P_2=P_3\)，故选 D.
\end{solution}

\begin{problem}
如图，小圆圈表示网络的结点，结点之间的连线表示它们有网线相联. 连线标注的数字表示该段网线单位时间内可以通过的最大信息量. 现从结点 \(A\) 向结点 \(B\) 传递信息，信息可以分开沿不同的路线同时传递. 则单位时间内传递的最大信息量为（\quad）

\bitmapfigure[width=6.3cm]{img/2001/henan_guangdong/q12_fig1.png}

\choices
    {\(26\)}
    {\(24\)}
    {\(20\)}
    {\(19\)}
\end{problem}

\begin{answer}
D
\end{answer}

\begin{solution}
从 \(B\) 到 \(A\) 的网络分为上、下两部分. 上部从 \(B\) 发出的两条边容量之和为 \(3+4=7\)，所以通过上部的流量不超过 \(7\). 下部最后通向 \(A\) 的边容量为 \(12\)，所以通过下部的流量不超过 \(12\). 因而总流量满足

\[
F\leqslant7+12=19
\]

这个上界可以达到：上部令容量为 \(3\) 的边通过 \(3\)，容量为 \(4\) 的边通过 \(4\)，再由上部右侧结点向 \(A\) 传递 \(7\)；下部令 \(B\) 到下方中间结点的边通过 \(6\)，\(B\) 到最下方结点的边通过 \(6\)，两股信息分别沿容量为 \(6\) 和 \(8\) 的边汇合，再由下部右侧结点向 \(A\) 传递 \(12\). 上述各边流量均不超过相应容量，所以总流量 \(7+12=19\) 可以实现. 因此单位时间内传递的最大信息量为 \(19\)，故选 D.
\end{solution}

\section{填空题}
\begin{problem}
已知甲、乙两组各有 \(8\text{ 人}\)，现从每组抽取 \(4\text{ 人}\) 进行计算机知识竞赛，比赛人员的组成共有\fillinblank{} 种可能.
\end{problem}

\begin{answer}
4900
\end{answer}

\begin{solution}
从甲组 \(8\) 人中抽取 \(4\) 人有 \(\mathrm C_8^4\) 种方法，从乙组 \(8\) 人中抽取 \(4\) 人也有 \(\mathrm C_8^4\) 种方法. 两组抽取相互独立，所以组成比赛人员的方法数为

\[
\mathrm C_8^4\mathrm C_8^4=70^2=4900
\]
\end{solution}

\begin{problem}
双曲线 \(\frac{x^{2}}{9} - \frac{y^{2}}{16} = 1\) 的两个焦点为 \(F_{1}\)，\(F_{2}\)，点 \(P\) 在双曲线上，若 \(PF_{1} \perp PF_{2}\)，则点 \(P\) 到 \(x\) 轴的距离为\fillinblank{}.
\end{problem}

\begin{answer}
\(\frac{16}{5}\)
\end{answer}

\begin{solution}
双曲线的实半轴、虚半轴分别为 \(3\)，\(4\)，所以焦半距为

\[
c=\sqrt{3^2+4^2}=5
\]

取两个焦点为 \(F_1(-5,0)\)，\(F_2(5,0)\)，设 \(P(x,y)\). 由 \(PF_1\perp PF_2\)，有

\[
\overrightarrow{PF_1}\cdot\overrightarrow{PF_2}=(-5-x,-y)\cdot(5-x,-y)=x^2+y^2-25=0
\]

即 \(x^2=25-y^2\). 代入双曲线方程，得

\[
\frac{25-y^2}{9}-\frac{y^2}{16}=1
\]

两边乘以 \(144\)，化简得

\(400-16y^2-9y^2=144\)，\(qquad 25y^2=256\).

所以点 \(P\) 到 \(x\) 轴的距离为 \(|y|=\frac{16}{5}\).
\end{solution}

\begin{problem}
设 \(\{a_{n}\}\) 是公比为 \(q\) 的等比数列，\(S_{n}\) 是它的前 \(n\) 项和. 若 \(\{S_{n}\}\) 是等差数列，则 \(q =\)\fillinblank{}.
\end{problem}

\begin{answer}
1
\end{answer}

\begin{solution}
由于 \(\{S_n\}\) 是等差数列，对 \(n\geqslant2\) 有

\[
S_n-S_{n-1}=a_n
\]

所以数列 \(\{a_n\}_{n\geqslant2}\) 是常数数列，特别地 \(a_2=a_3\). 又 \(\{a_n\}\) 是等比数列，故

\[
a_1q=a_1q^2
\]

等比数列的首项不为零且公比 \(q\ne0\)，于是 \(q(1-q)=0\) 只能给出 \(q=1\).
\end{solution}

\begin{problem}
圆周上有 \(2n\) 个等分点（\(n > 1\)），以其中三个点为顶点的直角三角形的个数为\fillinblank{}.
\end{problem}

\begin{answer}
\(2n(n-1)\)
\end{answer}

\begin{solution}
圆内接三角形是直角三角形，当且仅当其一条边为圆的直径. \(2n\) 个等分点中，任取一个点，其对面的点唯一确定，所以直径共有 \(n\) 条. 对于每条直径，第三个顶点可以从其余 \(2n-2\) 个等分点中任取，得到 \(2n-2\) 个直角三角形.

不同直径得到的直角三角形不重复，因为直角三角形的斜边直径是唯一的. 因此直角三角形的个数为

\[
n(2n-2)=2n(n-1)
\]
\end{solution}

\section{解答题}
\begin{problem}
求函数 \(y = (\sin x + \cos x)^2 + 2\cos^2 x\) 的最小正周期.
\end{problem}

\begin{solution}
由二倍角公式，
\[
y=(\sin x+\cos x)^2+2\cos^2x=1+\sin 2x+1+\cos 2x=2+\sqrt{2}\sin\left(2x+\frac{\pi}{4}\right)
\]
因为 \(\sin\left(2x+\frac{\pi}{4}\right)\) 的最小正周期为 \(\pi\)，且前面的系数 \(\sqrt{2}\) 不为零，所以函数 \(y\) 的最小正周期为 \(\pi\).
\end{solution}

\begin{problem}
已知等差数列前三项为 \(a\)，\(4\)，\(3a\)，前 \(n\) 项的和为 \(S_{n}\)，\(S_{k} = 2550\).

\begin{enumerate}
\item 求 \(a\) 及 \(k\) 的值；
\item 求 \(\displaystyle\lim_{n\to\infty}\left(\frac{1}{S_{1}}+\frac{1}{S_{2}}+\cdots+\frac{1}{S_{n}}\right)\).
\end{enumerate}
\end{problem}

\begin{solution}
\begin{enumerate}
\item 因为 \(a\)，\(4\)，\(3a\) 是等差数列的前三项，所以相邻两项的差相等，即
\[
4-a=3a-4
\]
解得 \(a=2\)，公差为 \(4-2=2\)，因此 \(a_n=2+2(n-1)=2n\)，从而
\[
S_n=\frac{n(a_1+a_n)}{2}=\frac{n(2+2n)}{2}=n(n+1)
\]
由 \(S_k=2550\) 得
\[
k(k+1)=2550
\]
即 \(k^2+k-2550=0\). 因为 \(k>0\)，且 \(1+4\times2550=10201=101^2\)，所以
\[
k=\frac{-1+101}{2}=50
\]
故 \(a=2\)，\(k=50\).
\item 由 \(S_n=n(n+1)\)，得
\[
\sum_{i=1}^{n}\frac{1}{S_i}=\sum_{i=1}^{n}\frac{1}{i(i+1)}=\sum_{i=1}^{n}\left(\frac{1}{i}-\frac{1}{i+1}\right)=1-\frac{1}{n+1}
\]
因此
\[
\displaystyle\lim_{n\to\infty}\left(\frac{1}{S_1}+\frac{1}{S_2}+\cdots+\frac{1}{S_n}\right)=\lim_{n\to\infty}\left(1-\frac{1}{n+1}\right)=1
\]
\end{enumerate}
\end{solution}

\begin{problem}
如图，在底面是直角梯形的四棱锥 \(S - ABCD\) 中，\(\angle ABC = 90^{\circ}\)，\(SA \perp\) 面 \(ABCD\)，\(SA = AB = BC = 1\)，\(AD = \frac{1}{2}\).

\begin{enumerate}
\item 求四棱锥 \(S - ABCD\) 的体积；
\item 求面 \(SCD\) 与面 \(SBA\) 所成的二面角的正切值.
\end{enumerate}

\bitmapfigure[width=4.8cm]{img/2001/henan_guangdong/q19_fig1.png}
\end{problem}

\begin{solution}
\begin{enumerate}
\item 因为底面是直角梯形，且 \(AD\parallel BC\)，所以 \(AB\) 是梯形的高，从而底面面积为
\[
S_{ABCD}=\frac{AD+BC}{2}\cdot AB=\frac{\frac{1}{2}+1}{2}\cdot1=\frac{3}{4}
\]
又因为 \(SA\perp\) 面 \(ABCD\)，所以四棱锥的高为 \(SA=1\). 因此
\[
V_{S-ABCD}=\frac{1}{3}S_{ABCD}\cdot SA=\frac{1}{3}\cdot\frac{3}{4}\cdot1=\frac{1}{4}
\]
\item 建立空间直角坐标系，取
\(A(0,0,0)\)，\(B(0,1,0)\)，\(C(1,1,0)\)，\(D\left(\frac{1}{2},0,0\right)\)，\(S(0,0,1)\).
则平面 \(SBA\) 的一个法向量为 \(\bs{n}_1=(1,0,0)\). 又
\(\overrightarrow{SC}=(1,1,-1)\)，\(\overrightarrow{SD}=\left(\frac{1}{2},0,-1\right)\)，
所以平面 \(SCD\) 的一个法向量可取
\[
\bs{n}_2=2\left(\overrightarrow{SD}\times\overrightarrow{SC}\right)=(2,-1,1)
\]
设两平面所成的锐二面角为 \(\varphi\)，则
\[
\cos\varphi=\frac{|\bs{n}_1\cdot\bs{n}_2|}{|\bs{n}_1||\bs{n}_2|}=\frac{2}{\sqrt{6}}=\sqrt{\frac{2}{3}}
\]
于是
\[
\tan\varphi=\frac{\sqrt{1-\cos^2\varphi}}{\cos\varphi}=\frac{1}{\sqrt{2}}=\frac{\sqrt{2}}{2}
\]
\end{enumerate}
\end{solution}

\begin{problem}
设计一幅宣传画，要求画面面积为 \(4840 \mathrm{~cm}^2\)，画面的宽与高的比为 \(\lambda\)（\(\lambda < 1\)），画面的上、下各留 \(8 \mathrm{~cm}\) 空白，左、右各留 \(5 \mathrm{~cm}\) 空白. 怎样确定画面的高与宽尺寸，能使宣传画所用纸张面积最小？如果要求 \(\lambda \in \left[\frac{2}{3}, \frac{3}{4}\right]\)，那么 \(\lambda\) 为何值时，能使宣传画所用纸张面积最小？
\end{problem}

\begin{solution}
设画面的高为 \(h\)，宽为 \(w\)，其中 \(0<\lambda<1\). 由 \(wh=4840\)，\(\frac{w}{h}=\lambda\)，可得
\(h=\sqrt{\frac{4840}{\lambda}}\)，\(w=\sqrt{4840\lambda}\).
纸张的宽为 \(w+10\)，高为 \(h+16\)，所以纸张面积为
\[
P=(w+10)(h+16)=5000+16w+10h=5000+\sqrt{4840}\left(16\sqrt{\lambda}+\frac{10}{\sqrt{\lambda}}\right)
\]
对 \(\lambda\) 求导，得
\[
P'(\lambda)=\sqrt{4840}\left(\frac{8}{\sqrt{\lambda}}-\frac{5}{\lambda^{3/2}}\right)=\frac{\sqrt{4840}(8\lambda-5)}{\lambda^{3/2}}
\]
当 \(0<\lambda<\frac{5}{8}\) 时，\(P'(\lambda)<0\)；当 \(\frac{5}{8}<\lambda<1\) 时，\(P'(\lambda)>0\). 所以纸张面积在 \(\lambda=\frac{5}{8}\) 时取得最小值. 此时
\(h=\sqrt{\frac{4840}{5/8}}=88\)，\(w=\frac{5}{8}\cdot88=55\).
即画面的高为 \(88\mathrm{~cm}\)，宽为 \(55\mathrm{~cm}\).

若 \(\lambda\in\left[\frac{2}{3},\frac{3}{4}\right]\)，则整个区间都在 \(\left(\frac{5}{8},1\right)\) 内，故 \(P'(\lambda)>0\)，纸张面积随 \(\lambda\) 增大而增大，最小值在左端点取得，即
\[
\lambda=\frac{2}{3}
\]
\end{solution}

\begin{problem}
已知椭圆 \(\frac{x^{2}}{2} + y^{2} = 1\) 的右准线 \(l\) 与 \(x\) 轴相交于点 \(E\)，过椭圆右焦点 \(F\) 的直线与椭圆相交于 \(A\)，\(B\) 两点，点 \(C\) 在右准线 \(l\) 上，且 \(BC \parallel x\) 轴，求证：直线 \(AC\) 经过线段 \(EF\) 的中点.
\end{problem}

\begin{solution}
由椭圆方程可知，长半轴 \(a_0=\sqrt{2}\)，短半轴 \(b_0=1\)，焦半距 \(c=\sqrt{a_0^2-b_0^2}=1\). 因此右焦点为 \(F(1,0)\)，离心率为 \(e=\frac{1}{\sqrt{2}}\)，右准线方程为
\[
x=\frac{a_0}{e}=2
\]
所以 \(E=(2,0)\)，线段 \(EF\) 的中点为 \(M\left(\frac{3}{2},0\right)\).

先设直线 \(AB\) 不是竖直直线，写成 \(y=k(x-1)\). 设 \(A\)，\(B\) 的横坐标分别为 \(x_1\)，\(x_2\)，则将直线方程代入椭圆方程，得到
\[
\left(\frac{1}{2}+k^2\right)x^2-2k^2x+(k^2-1)=0
\]
由韦达定理，
\[
x_1+x_2=\frac{4k^2}{1+2k^2}
\]
并且
\[
x_1x_2=\frac{2(k^2-1)}{1+2k^2}
\]
从而
\[
2x_1x_2-3(x_1+x_2)+4=0
\]
又因为 \(A=(x_1,k(x_1-1))\)，\(B=(x_2,k(x_2-1))\)，且 \(BC\parallel x\) 轴、\(C\in l\)，所以 \(C=(2,k(x_2-1))\). 直线 \(AC\) 在横坐标 \(\frac{3}{2}\) 处的纵坐标为
\[
\frac{\frac{1}{2}k(x_1-1)+\left(\frac{3}{2}-x_1\right)k(x_2-1)}{2-x_1}=0
\]
其中分子乘以 \(2\) 后等于 \(k\bigl(3x_1+3x_2-2x_1x_2-4\bigr)=0\)，上式确实成立. 因此 \(M\in AC\).

若直线 \(AB\) 为竖直直线，则其方程为 \(x=1\)，两个交点为 \(\left(1,\frac{1}{\sqrt{2}}\right)\) 和 \(\left(1,-\frac{1}{\sqrt{2}}\right)\). 无论哪一点记为 \(B\)，由 \(C\in l\) 且 \(BC\parallel x\) 轴可知，\(C\) 的坐标为 \(\left(2,\pm\frac{1}{\sqrt{2}}\right)\)，另一交点为 \(A\). 于是 \(M\left(\frac{3}{2},0\right)\) 是线段 \(AC\) 的中点，仍有 \(M\in AC\). 所以直线 \(AC\) 经过线段 \(EF\) 的中点.
\end{solution}

\begin{problem}
设 \(f(x)\) 是定义在 \(\R\) 上的偶函数，其图象关于直线 \(x = 1\) 对称，对任意 \(x_{1}\)，\(x_{2} \in \left[0, \frac{1}{2}\right]\)，都有 \(f(x_{1} + x_{2}) = f(x_{1}) \cdot f(x_{2})\)，且 \(f(1) = a > 0\).

\begin{enumerate}
\item 求 \(f\left(\frac{1}{2}\right)\)，\(f\left(\frac{1}{4}\right)\)；
\item 证明：\(f(x)\) 是周期函数；
\item 记 \(a_{n} = f\left(2n + \frac{1}{2n}\right)\)，求 \(\displaystyle\lim_{n \to \infty}(\ln a_{n})\).
\end{enumerate}
\end{problem}

\begin{solution}
\begin{enumerate}
\item 取 \(x_1=x_2=\frac{1}{2}\)，得
\[
f(1)=f\left(\frac{1}{2}\right)^2
\]
所以 \(f\left(\frac{1}{2}\right)=\pm\sqrt{a}\). 又取 \(x_1=x_2=\frac{1}{4}\)，有
\[
f\left(\frac{1}{2}\right)=f\left(\frac{1}{4}\right)^2\geqslant0
\]
从而 \(f\left(\frac{1}{2}\right)=\sqrt{a}\). 再取 \(x_1=x_2=\frac{1}{8}\)，得 \(f\left(\frac{1}{4}\right)=f\left(\frac{1}{8}\right)^2\geqslant0\)，结合
\[
f\left(\frac{1}{4}\right)^2=f\left(\frac{1}{2}\right)=\sqrt{a}
\]
可得 \(f\left(\frac{1}{4}\right)=a^{1/4}\).
\item 图象关于直线 \(x=1\) 对称，故对任意 \(x\in\R\)，有 \(f(2-x)=f(x)\). 又因为 \(f\) 是偶函数，所以
\[
f(x+2)=f(2-(-x))=f(-x)=f(x)
\]
因此 \(f(x)\) 是周期函数，\(2\) 是它的一个周期.
\item 先由 \(f(1)=a>0\) 及 \(f(1)=f\left(\frac{1}{2}\right)^2\) 知道 \(f\left(\frac{1}{2}\right)\neq0\). 由题设关系
\[
f\left(\frac{1}{2}\right)=f(0)f\left(\frac{1}{2}\right)
\]
可知 \(f(0)=1\). 对任意正整数 \(n\)，令 \(u=\frac{1}{2n}\). 对 \(j=1\)，\(2\)，\(\ldots\)，\(n\)，有 \((j-1)u\in\left[0,\frac{1}{2}\right]\)，因此可以反复使用题设关系，得到
\[
f(ju)=f((j-1)u)f(u)
\]
从而
\[
f\left(\frac{1}{2}\right)=f(nu)=f(u)^n
\]
另一方面，由题设关系 \(f(u)=f\left(\frac{u}{2}\right)^2\geqslant0\). 结合 \(f\left(\frac{1}{2}\right)=\sqrt{a}>0\) 与上式，得 \(f(u)>0\)，再得
\[
f\left(\frac{1}{2n}\right)=f(u)=\sqrt[n]{\sqrt{a}}=a^{1/(2n)}
\]
由第（2）问知 \(f\) 的周期为 \(2\)，故
\[
a_n=f\left(2n+\frac{1}{2n}\right)=f\left(\frac{1}{2n}\right)=a^{1/(2n)}>0
\]
因此
\[
\lim_{n\to\infty}\ln a_n=\lim_{n\to\infty}\frac{\ln a}{2n}=0
\]
\end{enumerate}
\end{solution}
